\section{引言}

积雪覆盖范围和雪水储量描述积雪变化的不同方面。覆盖率衡量地表有雪的比例，雪水当量衡量积雪所含液态水的等效深度。相同的覆盖面积可以对应不同的雪水当量。在冬季变暖背景下，若只跟踪覆盖范围，便无法判断积雪保有的水量是否同步变化。区域聚合也可能掩盖不同纬度和经度上方向相反的变化。

本文分析同一再分析子集的 43 个冬季，检验覆盖面积与雪水当量的长期趋势和年际异常是否同步，检验这种关系是否随区域和纬度变化，并估计去除线性时间趋势后积雪指标与 2 米气温的关联。最后用积雪覆盖率等值线检验覆盖边界是否移动，以及推断对阈值、孤立斑块和有效经度的依赖程度。我们预期这些关系存在空间差异。各指标的趋势方向、斜率差异及边界位移均由数据估计。

\section{数据与质量控制}

\subsection{分析域与数据来源}

输入为随任务提供的 ERA5 冬季教学子集。ERA5 产品来源见文献~\cite{hersbach2020}。归档文件含月度主文件、1980 冬季样本和冬季备用场。趋势分析只使用月度主文件。月度主文件含 1979 年 12 月至 2022 年 2 月间的 129 个 DJF 月，纬度 $0^\circ$--$90^\circ$N、经度 $-180^\circ$--$179^\circ$，格距 $1^\circ$。原变量为积雪覆盖率（snow cover fraction，SCF，0--1）、雪水当量（snow water equivalent，SWE，mm）和 2 米气温（T2m，$^\circ$C）。气温直接沿用原单位。

原生雪变量的固定有效范围含 12\,467 个格点。数据包的固定分析掩膜进一步排除 623 个在全部月均被标记为至少 9\,999~mm SWE 的模式固定冰格点，留下 11\,844 个格点，总面积 $98.289\times 10^6~\mathrm{km^2}$（图~\ref{fig:domain}）。权重来自包内球面格点面积 \texttt{cell\_area}。掩膜反映输入产品的固定可用范围，格点面积按完整格点计。三变量在该域的 43 个冬季均完整，完整时间覆盖敏感性分析与主分析同域。所有地图空白处表示掩膜外或缺测。研究未引入外部气象、地形或地理边界数据。

\fig{fig01_domain}{fig:domain}{\textbf{北半球冬季积雪研究的固定分析域。} 阴影为随数据提供的分析掩膜（11\,844 个格点，$1^\circ$ 网格）。线框为欧洲、亚洲和北美洲经纬度代表窗口，覆盖同一套 1980--2022 年 43 个 DJF 冬季的 SCF、SWE 和 T2m。窗口用于近似区域比较。}

\subsection{DJF 聚合与备用场核对}

冬季按 1--2 月所在年份命名。对每格点的 12、1、2 月值按实际天数加权，闰年 2 月取 29 天。只有三个月均有效才生成冬季值。分析前检查 SCF 范围、SWE 非负性以及月度 SWE 是否达到 9\,999~mm。掩膜外即使存在有限 SWE 或气温值也不纳入空间统计。

双精度冬季场与备用场在全网格及分析域内缺测位置均一致。最大绝对差依次为 $9.30\times10^{-8}$（SCF 比例）、$1.01\times10^{-3}$~mm（SWE）和 $4.32\times10^{-6}\,^\circ$C（T2m）。按月值乘权重及求和使用单精度、除以季节天数使用双精度，可逐格复现备用场。后续分析用双精度聚合。两份冬季场源于同一月度子集，备用场核验检验的是聚合过程。

\section{研究方法}

\subsection{空间聚合与储量口径}

设 $A_i$ 为格点面积，$M_i$ 为固定分析域或其子域的二元掩膜，$f_{it}$、$s_{it}$、$T_{it}$ 分别为冬季 SCF、SWE 和 T2m。在有效格点集合上计算
\begin{align}
  \mathrm{SCA}_t &= 10^{-6}\sum_i A_i M_i f_{it}\quad(\mathrm{km^2}),\\
  \mathrm{SCR}_t &= 100\,\frac{\sum_i A_i M_i f_{it}}{\sum_i A_i M_i}\quad(\%),\\
  \overline{\mathrm{SWE}}_t &= \frac{\sum_i A_i M_i s_{it}}{\sum_i A_i M_i}, \qquad
  \overline{T}_t = \frac{\sum_i A_i M_i T_{it}}{\sum_i A_i M_i}.
\end{align}
SCA 为积雪覆盖面积，SCR 为分析域内面积加权积雪覆盖率，$\overline{\mathrm{SWE}}$ 和 $\overline T$ 的单位分别为 mm 和 $^\circ$C。长期变化以百分点（pp）、mm 或 $^\circ$C 每十年表示。SCR 的 1 pp 对应 SCF 比例的 0.01。

提供的元数据未说明 SWE 是全格点平均，还是仅在有雪面积上的条件平均。假设 A 对应全格点平均，假设 B 对应有雪面积条件平均。两种冬季储量情景均用聚合后的冬季场计算
\begin{equation}
 V_{A,t} = 10^{-12}\sum_i A_iM_i s_{it}, \qquad
 V_{B,t} = 10^{-12}\sum_i A_iM_i f_{it}s_{it}\quad(\mathrm{km^3}).
\end{equation}
其中 mm 到 m 换算乘 $10^{-3}$，$\mathrm{m^3}$ 到 $\mathrm{km^3}$ 再乘 $10^{-9}$。在固定域内 $V_A$ 与 $\overline{\mathrm{SWE}}$ 线性等价。$V_B$ 为冬季均值的乘积，对应题设的季节尺度估计，月内及月际 SCF 与 SWE 的协变未纳入该估计。为量化 SWE 极高值的影响，另在全时段内始终排除任何一冬季 SWE 曾超过 5\,000、2\,000 或 1\,000~mm 的格点，逐口径重新估计基线和趋势。此分析用于衡量高值敏感性，保留完整固定掩膜为主结果。

\subsection{趋势、区域比较与覆盖--储量同步性}

按实际冬季年份拟合 $Y_t=\alpha+\beta(t-\bar t)+\epsilon_t$，报告 $10\beta$。区间与主检验采用 Bartlett 核、最大滞后 3 个冬季的 Newey--West HAC 协方差~\cite{newey1987}。源数据同时保存常规 OLS 检验、Mann--Kendall（MK）检验和 Sen 斜率~\cite{mann1945,sen1968}。MK 检验用于比较单调趋势，Sen 斜率为所有时间点对的斜率中位数。格点主图采用逐变量 Benjamini--Hochberg 校正后的 HAC $q<0.05$~\cite{bh1995}。面积比例以有效面积为分母，格点比例以有效格点数为分母。按共享网格边定义空间邻接，经度首尾相接。通过 999 次随机置换斜率位置评估全局 Moran $I$。相邻格点具有相关性，显著格点数不能视为独立样本数。

文献记录 ERA5 于 2004 年 1 月引入卫星积雪资料同化，雪变量可能出现水平突变~\cite{mortimer2020,urraca2023}。作为敏感性检验，在全域趋势模型中加入水平项。该项在 DJF 1980--2003 年为 0，在 2004 冬季为 $60/91$（1--2 月在冬季中所占天数），自 2005 年起为 1。另将 2004、2005、2006 冬季作为整季断点，并拟合断点两侧的分段趋势。全域温度回归和主阈值边界年度趋势另用同一水平项复算。由于阈值边界是非线性统计量，另剔除混合期 DJF 2004，以 2005 年为断点比较各窗口和阈值。

预设区域窗口为欧洲 $35$--$72^\circ$N、$15^\circ$W--$60^\circ$E，亚洲 $5$--$85^\circ$N、$60^\circ$E--$180^\circ$，北美洲 $15$--$85^\circ$N、$170$--$50^\circ$W。经度左闭右开，纬度含端点。区域窗口与固定分析掩膜相交后互不重叠。矩形边缘不能精确归属俄罗斯、土耳其、格陵兰或北非和中亚边缘。纬度带按 $[0,30)$、$[30,50)$、$[50,65)$、$[65,90]^\circ$N 划分。通过调整窗口边界检验趋势方向与区域排序。

覆盖与雪水指标的同步性从三方面判断。第一，比较 SCA 与 $\overline{\mathrm{SWE}}$ 或储量情景的原始及去趋势 Pearson 相关。第二，计算 $D_t=z(\mathrm{SCA}_t)-z(\mathrm{SWE\ metric}_t)$。$D_t>0$ 表示该冬季覆盖面积相对自身基线的标准化异常高于雪水指标。第三，对 $D_t$ 的斜率及区域间斜率差采用同冬季配对的残差区块自助法，主区块长 5 个冬季，并比较 3 和 7 个冬季。配对重抽样保留同一冬季多序列的共同扰动，区块保留短期时间依赖~\cite{kunsch1989}。原始斜率具有不同量纲，因此直接比较使用标准差每十年。

\subsection{温度关联与跨冬季滞后}

先对各区域及纬度带的 SCF、SWE 和 T2m 分别移除线性时间趋势，再以温度残差为自变量拟合积雪残差。去趋势残差的 Pearson 相关等于控制年份的线性偏相关。报告残差斜率、HAC 区间、$R^2$、样本量、残差诊断及逐冬季删除后的斜率范围。用共同冬季的配对推断检验纬度带或区域间温度斜率差，避免把一组显著而另一组不显著当作组间差异。

设 $X_t$ 为 SCF 或 SWE 的去趋势残差，$\widetilde T_t$ 为温度残差。计算 $r_k=\mathrm{corr}(X_t,\widetilde T_{t-k})$，$k=-2,-1,0,1,2$。正滞后表示气温领先，负滞后表示积雪领先。有效配对数为 $43-|k|$。分析只使用跨冬季配对。对同一滞后检验族校正多重比较。

\subsection{积雪覆盖阈值边界}

阈值雪线是 SCF 等值线定义的二维地理积雪边界，纬度指标刻画积雪覆盖南缘的位置。其单位为纬度度数，海拔不参与定义。主阈值 $\theta=0.50$，敏感性阈值为 0.25 与 0.75。分别对 1980--1989 年和 2013--2022 年的 10 个 DJF 冬季求平均 SCF，绘制早晚等值线和差值场。等值线只在有效分析网格内提取，掩膜外不填零。经度首尾按周期网格连接。

年度纬度指标从相邻有效网格间由南向北的低于阈值到达到阈值的交叉处作线性插值。每个经度取最南一处合格交叉。主规则要求阈值以上斑块至少含 9 个四邻接格点，交叉处北侧至少连续 2 个达到阈值的有效格点。敏感性规则分别使用 1 格点与 1 个连续格点、25 格点与 3 个连续格点。缺少合格交叉的经度记为缺测。报告经度中位数、四分位距和有效列数，并在全部冬季始终有交叉的经度上复算年度趋势。早晚位移是在两期均有交叉的经度上先求晚减早，再取经度中位数。空间区块自助区间反映给定两个时期下的经度抽样变异。年度趋势按每种斑块规则和经度支持下的 12 个阈值与区域组合进行 BH 校正。
